\documentclass{article}

\input{preamble.tex}

\title{Higher algebra in characteristic 2}
\author{Jacob Saunders}

\begin{document}

\maketitle

Everything in the talk is implicitly derived.
For example, (co)limits are really homotopy (co)limits when this makes sense, smash products are derived smash products, etc.

The bulk of this talk is really purloined from \cite{szymik}.

\section{The story, up to homotopy}

\begin{definition}
    A \textit{graded Eilenberg--MacLane spectrum} (of characteristic 2), or \textit{GEM} for short, is a spectrum of the form $\bigwedgesum_{\mathcal{I}} \Sigma^i \HF_2$ for some indexing multiset $\mathcal{I}$.
\end{definition}

Alternatively, you could think of these as the Eilenberg--MacLane spectra $HR$ of graded $\F_2$-algebras $R$.
These admit canonical ring structures up to homotopy, and can even be refined to $E_\infty$ rings.

The homotopy theory of these objects is algebraic, as is covered in \cite{boardman}.

\begin{lemma}[{\cite[Lemma 2.1]{boardman}}]
    Let $X$ be a GEM.
    Then the Hurewicz map $\pi_* X \to H_* X$ is injective, the embedded submodule given by the elements $x \in H_* X$ such that the dual Steenrod comodule action sends $x$ to $1 \otimes x$.
\end{lemma}

\begin{lemma}[{\cite[Lemma 2.3]{boardman}}]
    Let $X, Y$ be GEMs.
    Then (homotopy classes of) maps $f : X \to Y$ correspond to commutative squares
    % https://q.uiver.app/#q=WzAsNCxbMCwwLCJcXHBpXyogWCJdLFsxLDAsIlxccGlfKiBZIl0sWzAsMSwiSF8qIFgiXSxbMSwxLCJIXyogWSJdLFswLDEsImYiXSxbMCwyLCJoIiwyXSxbMSwzLCJoIl0sWzIsMywiZiIsMl1d
    \[\begin{tikzcd}
        {\pi_* X} & {\pi_* Y} \\
        {H_* X} & {H_* Y}
        \arrow["f", from=1-1, to=1-2]
        \arrow["h"', from=1-1, to=2-1]
        \arrow["h", from=1-2, to=2-2]
        \arrow["f"', from=2-1, to=2-2]
    \end{tikzcd}\]
    where the top map is a module map, and the bottom is a map of $\A_*$-comodules.
\end{lemma}

This result can be used prove that there are no exotic homotopy ring spectrum structures on GEMs.

\begin{lemma}[{\cite[Theorem 1.1]{boardman}}]
    Let $X$ be a homotopy ring spectrum whose underlying object is a GEM.
    Then there is an isomorphism of ring spectra $X \simeq H\pi_*X$.
\end{lemma}

It is feasible that there are homotopy ring spectra with $\pi_* X$ a graded $\F_2$-algebra but which are not additively GEMs.
If we impose a commutativity condition, however, then this is not the case.

\begin{theorem}[{\cite{pazhitnov-rudyak}, \cite[Theorem 1.1]{wurgler}}]
    If $X$ is a homotopy commutative ring spectrum of characteristic 2, then $X$ is additively a GEM.
\end{theorem}

Here, $X$ having characteristic 2 simply means that $\pi_0 X$ has characteristic 2.
Commutativity here is required -- Wurgler gives some examples (e.g.\ spectra representing bordism with singularities) for which this doesn't hold.

\section{(Weakly) initial $E_n$ rings of characteristic 2}
The initial rings $Z/p$ of characteristic $p$ are pretty fundamental, and it's natural to ask if there are analogues of these in higher algebra.
The ring $\Z/2$ can be thought of as the cofiber $\Z \xrightarrow{\cdot 2} \Z$, and so the obvious candidate in our situation is the cofiber $S^0 \xrightarrow{\cdot 2} S^0 \to S^0/2$.
This is given by the mod 2 Moore spectrum, but this does not admit a homotopy ring structure!

Another way to present $\Z/2$ is to take a pushout
% https://q.uiver.app/#q=WzAsNCxbMCwwLCJcXFpbeF0iXSxbMSwwLCJcXFoiXSxbMCwxLCJcXFoiXSxbMSwxLCJcXFovMiJdLFswLDEsInggXFxtYXBzdG8gMiJdLFswLDIsInggXFxtYXBzdG8gMCIsMyx7Im9mZnNldCI6Mywic3R5bGUiOnsiYm9keSI6eyJuYW1lIjoibm9uZSJ9LCJoZWFkIjp7Im5hbWUiOiJub25lIn19fV0sWzAsMl0sWzEsM10sWzIsM11d
\[\begin{tikzcd}
	{\Z[x]} & \Z \\
	\Z & {\Z/2}
	\arrow["{x \mapsto 2}", from=1-1, to=1-2]
	\arrow["{x \mapsto 0}"{marking, allow upside down}, shift right=3, draw=none, from=1-1, to=2-1]
	\arrow[from=1-1, to=2-1]
	\arrow[from=1-2, to=2-2]
	\arrow[from=2-1, to=2-2]
\end{tikzcd}\]
in the category of commutative rings, which is given by the relative tensor product $\Z \otimes_{Z[x]} \Z$.
In this context, we are guaranteed a ring structure.

To guarantee the existence of such objects, we work instead inside the category of $E_\infty$-algebras.

\begin{definition}[{\cite[Definition 2.1]{szymik}}]
    The \textit{versal $E_\infty$ ring spectrum of characteristic 2} is given by the pushout
    % https://q.uiver.app/#q=WzAsNCxbMCwwLCJcXFAoU14wKSJdLFsxLDAsIlxcUyJdLFswLDEsIlxcUyJdLFsxLDEsIlxcUyBcXG1vZG1vZCAyIl0sWzAsMiwiXFxldiAwIiwyXSxbMCwxLCJcXGV2IDIiXSxbMiwzXSxbMSwzXSxbMywwLCIiLDEseyJzdHlsZSI6eyJuYW1lIjoiY29ybmVyIn19XV0=
    \[\begin{tikzcd}
        {\P(S^0)} & \S \\
        \S & {\S \modmod 2}
        \arrow["{\ev 2}", from=1-1, to=1-2]
        \arrow["{\ev 0}"', from=1-1, to=2-1]
        \arrow[from=1-2, to=2-2]
        \arrow[from=2-1, to=2-2]
        \arrow["\lrcorner"{anchor=center, pos=0.125, rotate=180}, draw=none, from=2-2, to=1-1]
    \end{tikzcd}\]
    in the category of $E_\infty$ ring spectra.
\end{definition}
Here, $\P(-)$ denotes the free functor $\Spectra \to \Alg_{E_\infty}(\Spectra)$, and the maps $\ev(-)$ arise from the adjunction equivalence.
This construction can be thought of as freely adding a nullhomotopy of 2 to the sphere, and so if an $E_\infty$ ring $X$ has characteristic 2 (i.e.\ has a nullhomotopy of 2), then it recieves a map from $\S \modmod 2$.In fact, from its description as a pushout, we find
% https://q.uiver.app/#q=WzAsNCxbMCwwLCJcXEhvbV97RV9cXGluZnR5fShcXFAoU14wKSwgWCkiXSxbMSwwLCJcXEhvbV97RV9cXGluZnR5fShcXFMsIFgpIl0sWzAsMSwiXFxIb21fe0VfXFxpbmZ0eX0oXFxTLCBYKSJdLFsxLDEsIlxcSG9tX3tFX1xcaW5mdHl9KFxcUyBcXG1vZG1vZCAyLCBYKSJdLFszLDJdLFszLDFdLFsyLDBdLFsxLDBdLFszLDAsIiIsMSx7InN0eWxlIjp7Im5hbWUiOiJjb3JuZXIifX1dXQ==
\[\begin{tikzcd}
	{\Hom_{E_\infty}(\P(S^0), X)} & {\Hom_{E_\infty}(\S, X)} \\
	{\Hom_{E_\infty}(\S, X)} & {\Hom_{E_\infty}(\S \modmod 2, X)}
	\arrow[from=1-2, to=1-1]
	\arrow[from=2-1, to=1-1]
	\arrow["\lrcorner"{anchor=center, pos=0.125, rotate=225}, draw=none, from=2-2, to=1-1]
	\arrow[from=2-2, to=1-2]
	\arrow[from=2-2, to=2-1]
\end{tikzcd}\]
and using that $\Hom_{E_\infty}(\S, X) \simeq \ast$ because $\S$ is initial, and using the adjunction equivalence $\Hom_{E_\infty}(\P(S^0), X) \simeq \Hom_{\Spectra}(S^0, X) \simeq \Omega^\infty X$, the pullback is equivalent to
% https://q.uiver.app/#q=WzAsNCxbMCwwLCJcXE9tZWdhXlxcaW5mdHkgWCJdLFsxLDAsIlxcYXN0Il0sWzAsMSwiXFxhc3QiXSxbMSwxLCJcXEhvbV97RV9cXGluZnR5fShcXFMgXFxtb2Rtb2QgMiwgWCkiXSxbMywyXSxbMywxXSxbMiwwLCIwIl0sWzEsMCwiMiIsMl0sWzMsMCwiIiwxLHsic3R5bGUiOnsibmFtZSI6ImNvcm5lciJ9fV1d
\[\begin{tikzcd}
	{\Omega^\infty X} & \ast \\
	\ast & {\Hom_{E_\infty}(\S \modmod 2, X)}
	\arrow["2"', from=1-2, to=1-1]
	\arrow["0", from=2-1, to=1-1]
	\arrow["\lrcorner"{anchor=center, pos=0.125, rotate=180}, draw=none, from=2-2, to=1-1]
	\arrow[from=2-2, to=1-2]
	\arrow[from=2-2, to=2-1]
\end{tikzcd}\]
so this is empty is 2 is not nullhomotopic, and otherwise is equivalent to $\Omega(\Omega^{\infty} X)$.
Therefore maps out of these are not even unique up to homotopy if $\pi_1(X)$ is nontrivial.

This whole story admits a generalisation to $E_n$-algebras too, as developed in \cite{simple-universal-property}.
We denote the corresponding $E_n$ pushout by $\S \modmod_n 2$.

\subsection*{Equivalent methods of obtaining $\S \modmod_n 2$}
In the $E_\infty$ case, the pushout is a relative smash product, so there is an equivalence
\[
    \S \modmod 2 \simeq \S \smashprod_{\P(S^0)} \S
\]
of $E_\infty$ ring spectra.
Note that the $\P(S^0)$-module structures on the left and right smash factors are different.
In the $E_n$ case, work of Hill--Lawson \cite[Corollary 7.4]{hill-lawson} allows us to write some pushouts as relative smash products.
This gives us an equivalence
\[
    \S \modmod_n 2 \simeq \S \smashprod_{\P_{n + 1}(S^0)} \S.
\]

These can also be obtained as Thom spectra of the map $\Omega^n \Sigma^n S^1 \to BO$ induced by the non-trivial homotopy class $S^1 \to BO$.
This gives us the homology
\[
    H_*\S \modmod_n 2 \cong H_* \Omega^n\Sigma^nS^1 \cong \F_2[a, Q_1^{i_1} \cdots Q_{n - 1}^{i_{n - 1}}a],
\]
which had previously been computed by Cohen. % TODO would be good to cite that.

Since all these things are homotopy commutative and have characteristic 2, you can determine the homotopy by `dividing by' the dual Steenrod algebra.
Since
\[
    H_* \S \modmod_2 2 \cong \A_*,
\]
this means $\pi_* \S \modmod_2 2 \cong \F_2$, meaning $\S \modmod_2 2 \simeq \HF_2$.

\section{Applications to $E_n$ maps}
Because $\HF_2$ is $E_\infty$ and has homotopy concentrated in degree 0, there is an essentially unique map $\S \modmod 2 \to \HF_2$.

\begin{theorem}[{\cite[Proposition 5.1]{szymik}}]
    There is no $E_\infty$ map $\HF_2 \to \S \modmod 2$.    
\end{theorem}

\begin{proof}
    Any such map postcomposes with $\S \modmod 2 \to \HF_2$ to give the identity (up to homotopy) by our classification of $E_n$ endomorphisms of versal algebras, so sends $\xi_1$ to $a$.
    Furthermore, the map must commute with the Dyer--Lashof operations, predicting that $Q_2 a = Q_2 f(\xi_1) = f(Q_2\xi_1) = f(\xi^4) = a^4$.
    Since this relation does not hold in $\S \modmod 2$, we derive a contradiction.
\end{proof}

\begin{definition}
    The \textit{unoriented bordism spectrum}, denoted $MO$, is the Thom spectrum of the delooped $j$-homomorphism $BO \to B\text{GL}_1\S$.
\end{definition}

A result of Thom's original paper is that $\pi_* MO$ is the unoriented bordism ring, permitting him to compute that
\[
    \pi_* MO \cong \F_2[x_i : i \neq 2^j - 1].
\]
Knowing that $MO$ is $E_\infty$ and that the bordism ring has characteristic 2 allows us to use the results of the first part to show that $MO$ is a GEM.

If this is the whole story, then one might expect that there is an $E_\infty$ equivalence $H\pi_* MO \simeq MO$.
This is not the case!

\begin{theorem}[Helen Gilmour's Thesis]
    There is no $E_4$ map $\HF_2 \to MO$.
\end{theorem}

The existence of a map is obstructed by a Dyer--Lashof operation (it must be a $Q_3$).
These were computed in $\HF_2$ in the $H_\infty$ book, and in $BO$ (and therefore $MO$) by Priddy/Kochman.

To try and recover from this, one might expect that if $\pi_* MO$ is a polynomial ring, $MO$ might be `polynomial' over $\S \modmod 2$, but this is not the case!

\begin{theorem}[{\cite[Corollary 5.11]{szymik}}]
    There is no spectrum $X$ such that there exists an equivalence
    \[
        MO \simeq \S \modmod 2 \smashprod \P(X)
    \]
    of $E_\infty$ ring spectra.
\end{theorem}

\begin{proof}
    Otherwise there would be a retraction
    \[
        \S \modmod 2 \to MO \simeq \S \modmod 2 \smashprod \P(X) \to \S \modmod 2 \smashprod \P(*) \simeq \S \modmod 2,
    \]
    but $H_* MO$ is too small to support this.
    Actually, we can use another Dyer--Lashof operation argument: the image of $Q_2 a$ is $x_1^4$, so the map $\S \modmod 2 \to MO$ is not a homology injection.
\end{proof}

This raises the question -- what does $MO$ look like as an $E_\infty$ ring spectrum?

\section{Highly structured models of $MO$}
It turns out that the combinatorics of $H_* BO$ as an algebra over the $E_2$ Dyer--Lashof operations is sufficient to determine that
\[
    MO \simeq_{E_1} \HF_2 \smashprod \bigsmashprod_{i \neq 2^j - 1} \P_1(S^i),
\]
and
\[
    MO \simeq_{E_2} \HF_2 \smashprod \bigsmashprod_{i \in \mathbb{Z}_{> 0}} \P_2(S^{2i}),
\]
so Thom's picture is accurate on the level of $E_1$ and $E_2$ ring spectra.
Note that the smash product is not the coproduct of $E_n$ rings for $n \neq \infty$.

We know (from the earlier result about the non-existence of an $E_4$ map $\HF_2 \to MO$) that this won't be the case for $E_4$, but what about $E_3$?

\subsection*{A quick aside on $H_\infty$ and $E_\infty$}
The calculations I'm going to talk about only involve the $H_n$ structure of $MO$ (the Dyer--Lashof operations), but turn out to give $E_n$ models of $MO$.
Noel and Lawson both construct instances of $H_\infty$ ring spectra that do not refine to $E_\infty$ ring spectra, and Johnson--Noel construct an obstruction theory that deals with this.
Computing the obstruction groups is hard.
It's possible that in our cases of interest, the obstruction theory allows for our $H_\infty$ maps to be refined in a unique or canonical way, but the calculations I am going to talk about are somehow elementary in comparison.

There are no obstructions to the existence of an $E_3$ map $\HF_2 \to MO$ on the level of homology, but it remains to try and construct one.
If we can do this, the combinatorics of $H_* BO$ immediately give an equivalence
\[
    MO \simeq_{E_3} \HF_2 \smashprod \P_3(S^2) \smashprod \bigsmashprod_{i \in \mathbb{Z}_{> 0}} \P_3(S^{4i}).
\]

\subsection*{An $E_3$ model of $\HF_2$}
I still have things to check here, so take this last part of the talk with a pinch of salt.

Our aim is to equip $\HF_2$ with an $E_3$ universal property, so that construction of an $E_3$ map $\HF_2 \to MO$ becomes trivial.
The ideal thing would be to approximate $\HF_2$ using $E_3$ pushouts, iteratively adding and killing classes in the relevant degrees to get successively better approximations.
However, computing non-$E_3$ invariants of $E_3$ pushouts is difficult.
Results of Hill--Lawson from earlier, however, tell us that some $E_3$ pushouts are equivalent to relative smash products, and invariants of these can be computed using the Kunneth spectral sequence.

The pushouts that can be computed as relative smash products come from \textit{actions}, classified by homotopy classes $S^i \to Z_n(X)$, where $Z_n$ denotes the \textit{$E_n$ centre} of $X$.
We can iteratively construct the relevant homotopy classes using universal properties of more structured variants of $X$.

\bibliographystyle{alpha}
\bibliography{refs}

\end{document}